\section{Background}
\label{sec:background}

\begin{figure*}[t]
  \centering
  \begin{subfigure}{0.49\textwidth}
    \centering
    \includegraphics[width=\linewidth]{figs/para_dist/paper_Qwen3-4B_hidden_both_block_para_perp_sampled_layers.pdf}
    \caption{Qwen3-4B-Instruct-2507}
  \end{subfigure}\hfill
  \begin{subfigure}{0.49\textwidth}
    \centering
    \includegraphics[width=\linewidth]{figs/para_dist/paper_Qwen3-30B-A3B_hidden_both_block_para_perp_sampled_layers.pdf}
    \caption{Qwen3-30B-A3B}
  \end{subfigure}
  \caption{\textbf{Parallel components persist across depth.}
  For each accumulated attention-plus-MLP update, $\Dup{l}{t}$ is its
  projection along the incoming hidden state, whereas $\Duperp{l}{t}$ is
  orthogonal to that state. Curves show the ratio
  $\|\Dup{l}{t}\|/\|\Duperp{l}{t}\|$ at three sampled layers per model
  across generation steps; lines are means and bands span the observed range.
  Ratios above one indicate parallel dominance.}
  \label{fig:profiles}
\end{figure*}

\subsection{Parallel and Perpendicular Decomposition}
\label{sec:decomp}

Given a reference vector $\mathbf{z} \in \mathbb{R}^d$ and an update vector
$\Delta \in \mathbb{R}^d$ within the same representation space, $\Delta$ admits
a unique orthogonal decomposition into components parallel and perpendicular to $\mathbf{z}$:
\begin{align}
  \Delta_\parallel
    &= \frac{\Delta \cdot \mathbf{z}}{\|\mathbf{z}\|^2}\,\mathbf{z}
     = \alpha\,\mathbf{z},
     \qquad
     \alpha = \frac{\Delta \cdot \mathbf{z}}{\|\mathbf{z}\|^2},
  \label{eq:para} \\
  \Delta_\perp
    &= \Delta - \Delta_\parallel,
    \qquad \Delta_\perp \perp \mathbf{z}.
  \label{eq:perp}
\end{align}
The parallel component $\Delta_\parallel$ lies along the reference direction
$\mathbf{z}$, scaling its magnitude by $\alpha$, whereas $\Delta_\perp$ lies in
the $(d-1)$-dimensional orthogonal subspace, carrying the direction-changing
degrees of freedom.

We apply this decomposition across two primary sites: in
\textit{residual space}, $\Delta$ is the sub-layer update added to the incoming
state $\mathbf{z}$; in \textit{attention value space}, $\Delta$ is the
pre-projection aggregate relative to the token's own value vector $\mathbf{z}$.

\subsection{Transformer Residual Updates}
\label{sec:transformer}

In a transformer with hidden dimension $d$, each residual sub-layer adds an update to the current representation:
\begin{equation}
  \hl{l+1} = \hl{l} + \Dl{l}.
  \label{eq:residual}
\end{equation}
The update $\Dl{l}$ denotes the output of an individual sub-layer
(self-attention $\Attn$ or $\mathrm{MLP}$) or the net update of the full
transformer block. The reference state $\hz{l}{t}$ is taken as the
representation immediately preceding the addition: the sub-layer input for
individual modules, or the block input for the accumulated update. We write
$\Du{l}{t}$ for the corresponding residual contribution at token position
$t \in \{1,\ldots,T\}$.

Using the decomposition in Equation~\ref{eq:para}, the parallel component of a residual update is:
\begin{equation}
  \Dup{l}{t} = \alpha^l_t \hz{l}{t}, \qquad
  \alpha^l_t = \frac{\Du{l}{t}\cdot \hz{l}{t}}{\|\hz{l}{t}\|^2},
\end{equation}
where the scalar $\alpha^l_t$ measures the magnitude of the parallel update relative to the incoming state $\hz{l}{t}$. Substituting this decomposition into the residual addition yields
\begin{equation}
  \hz{l+1}{t}
  = (1+\alpha^l_t)\hz{l}{t} + \Duperp{l}{t}, \qquad
  \Duperp{l}{t}\perp \hz{l}{t},
\end{equation}
where the parallel component rescales the incoming representation by $(1+\alpha^l_t)$, while $\Duperp{l}{t}$ introduces orthogonal direction-changing information.

Figure~\ref{fig:profiles} reports the relative magnitude $\|\Dup{l}{t}\|/\|\Duperp{l}{t}\|$ across generation steps at three sampled layers for two Qwen3 models. This ratio quantifies how much of a residual update aligns with the existing representation versus redirects it. The persistence of substantial parallel components across depth motivates investigating when they affect model behavior and in which representation spaces those effects remain robust.

\begin{figure*}[t]
  \centering
  \includegraphics[width=0.99\textwidth]{figs/overview/parallel_updates_main_figure_analysis_v11_midline_arrows.pdf}
  \caption{\textbf{Geometric decomposition and component scaling.}
  Residual updates and attention value aggregates are decomposed into parallel
  and perpendicular components and scaled independently; parallel-only scaling
  can be expressed as an attention-diagonal change. Labels report raw $\Delta\mathrm{PPL}$ relative to the
  scale-one no-op (lower is better); curves are rescaled only for display.}
  \label{fig:overview}
\end{figure*}
